% =============================================================================
% Chapter 11: Appendix
% =============================================================================
\chapter{Appendix}
\label{ch:appendix}

\section{Appendix A: Complete RTL File Inventory}

\begin{table}[H]
\centering
\caption{RTL Source Files}
\label{tab:rtl_files}
\begin{tabularx}{\textwidth}{@{}lrX@{}}
\toprule
\textbf{File} & \textbf{Lines} & \textbf{Description} \\
\midrule
\texttt{apb\_spi\_master.sv}       & 378  & Top-level module with interrupt logic and sub-module instantiation. \\
\texttt{spi\_master\_apb\_if.sv}   & 220  & APB register interface (address decode, read mux, write logic). \\
\texttt{spi\_master\_controller.sv}& 561  & Controller FSM with state machine, data mux, and CS logic. \\
\texttt{spi\_master\_clkgen.sv}    & 82   & SPI clock generator with programmable divider. \\
\texttt{spi\_master\_tx.sv}        & 130  & TX shift register (standard and quad mode). \\
\texttt{spi\_master\_rx.sv}        & 140  & RX shift register (standard and quad mode). \\
\texttt{spi\_master\_fifo.sv}      & 140  & Parameterized FIFO with circular buffer. \\
\texttt{spi\_master\_tb\_v2.sv}    & 568  & Testbench v2 with slave BFM and 14 testcases. \\
\midrule
\textbf{Total}                     & \textbf{2219} & \\
\bottomrule
\end{tabularx}
\end{table}

\section{Appendix B: Simulation Commands}

\subsection{Compiling and Running with Xcelium}

\begin{lstlisting}[language=bash,numbers=none]
# Compile all RTL and testbench files
xmvlog -sv spi_master_fifo.sv \
           spi_master_clkgen.sv \
           spi_master_tx.sv \
           spi_master_rx.sv \
           spi_master_controller.sv \
           spi_master_apb_if.sv \
           apb_spi_master.sv \
           spi_master_tb_v2.sv

# Elaborate the testbench top module
xmelab tb_apb_spi_master

# Run the simulation
xmsim tb_apb_spi_master
\end{lstlisting}

\subsection{Using a File List}

Create a file \texttt{spi\_file.f} containing:
\begin{lstlisting}[numbers=none]
spi_master/spi_master_fifo.sv
spi_master/spi_master_clkgen.sv
spi_master/spi_master_tx.sv
spi_master/spi_master_rx.sv
spi_master/spi_master_controller.sv
spi_master/spi_master_apb_if.sv
spi_master/apb_spi_master.sv
spi_master/spi_master_tb_v2.sv
\end{lstlisting}

Then compile with:
\begin{lstlisting}[language=bash,numbers=none]
xmvlog -sv -f spi_file.f
xmelab tb_apb_spi_master
xmsim tb_apb_spi_master
\end{lstlisting}

\subsection{Viewing Waveforms in SimVision}

\begin{lstlisting}[language=bash,numbers=none]
simvision waves_v2.trn
\end{lstlisting}

\section{Appendix C: Register Quick Reference Card}

\begin{table}[H]
\centering
\caption{Register Quick Reference}
\label{tab:reg_quick}
\small
\begin{tabular}{@{}lll@{}}
\toprule
\textbf{Offset} & \textbf{Name} & \textbf{Key Fields} \\
\midrule
0x00 & STATUS  & [0]=rd, [1]=wr, [2]=qrd, [3]=qwr, [4]=swrst, [11:8]=csreg \\
0x04 & CLKDIV  & [7:0]=divider \\
0x08 & SPICMD  & [31:0]=command value \\
0x0C & SPIADR  & [31:0]=address value \\
0x10 & SPILEN  & [5:0]=cmd\_len, [13:8]=addr\_len, [31:16]=data\_len \\
0x14 & SPIDUM  & [15:0]=dummy\_rd, [31:16]=dummy\_wr \\
0x18 & TXFIFO  & [31:0]=write data (push) \\
0x20 & RXFIFO  & [31:0]=read data (pop) \\
0x24 & INTCFG  & [30]=cnt\_en, [31]=int\_en, thresholds and counters \\
0x28 & INTSTA  & Read to re-arm interrupt FSM \\
\bottomrule
\end{tabular}
\end{table}

\section{Appendix D: FSM State Encoding Quick Reference}

\begin{table}[H]
\centering
\caption{Controller FSM State Encoding}
\begin{tabular}{@{}cl@{}}
\toprule
\textbf{Value} & \textbf{State} \\
\midrule
0 & IDLE \\
1 & CMD \\
2 & ADDR \\
3 & MODE \\
4 & DUMMY \\
5 & DATA\_TX \\
6 & DATA\_RX \\
7 & WAIT\_EDGE \\
\bottomrule
\end{tabular}
\end{table}

\section{Appendix E: Testbench Task Quick Reference}

\begin{table}[H]
\centering
\caption{Testbench Task Summary}
\label{tab:task_ref}
\small
\begin{tabularx}{\textwidth}{@{}lX@{}}
\toprule
\textbf{Task/Function} & \textbf{Description} \\
\midrule
\texttt{reg\_addr(idx)}        & Converts 4-bit register index to 12-bit APB address. \\
\texttt{apb\_write(addr,data)} & Performs a 2-cycle APB write transaction. \\
\texttt{apb\_read(addr,rdata)} & Performs a 2-cycle APB read transaction with correct FIFO timing. \\
\texttt{check(msg,got,exp,mask)} & Masked comparison with pass/fail counting and formatted output. \\
\texttt{cfg\_spi(...)}         & Writes all 5 SPI configuration registers in one call. \\
\texttt{trigger(op,cs)}       & Writes STATUS register to start an SPI transaction. \\
\texttt{push\_tx(data)}       & Pushes a 32-bit word into the TX FIFO. \\
\texttt{pop\_rx(data)}        & Pops a 32-bit word from the RX FIFO. \\
\texttt{wait\_eot(timeout)}   & Waits for the sticky EOT flag with configurable timeout. \\
\texttt{clear\_eot\_flag}     & Resets the sticky EOT flag using force/release. \\
\bottomrule
\end{tabularx}
\end{table}

\section{Appendix F: Known Limitations and FIXME Notes}

The following FIXME comments exist in the RTL source code:

\begin{enumerate}[leftmargin=2em]
  \item \texttt{spi\_master\_apb\_if.sv}, line 284: \texttt{spi\_data\_tx\_ready} is marked as
        ``not used inside this module.'' The signal is used by the FIFO, not the APB interface.
        This is informational, not a bug.
  \item \texttt{spi\_master\_controller.sv}, line 31: \texttt{spi\_swrst} is marked as ``not used
        at all'' within the controller. The software reset only clears the FIFOs at the top level.
        The controller itself does not respond to software reset, which means issuing swrst during
        an active transaction may leave the FSM in an inconsistent state.
  \item \texttt{spi\_master\_controller.sv}, line 77: \texttt{do\_tx} is marked as ``NOT USED
        at all.'' This signal is set but never read in any conditional logic. It can be safely
        removed in a future cleanup.
\end{enumerate}

\section{Appendix G: Document Revision History}

\begin{table}[H]
\centering
\caption{Document Revision History}
\begin{tabularx}{\textwidth}{@{}cclX@{}}
\toprule
\textbf{Rev} & \textbf{Date} & \textbf{Author} & \textbf{Changes} \\
\midrule
1.0 & May 2026 & Verification Team & Initial release. 14 testcases, 28 checks, all passing. \\
\bottomrule
\end{tabularx}
\end{table}

\vfill
\begin{center}
\rule{0.5\textwidth}{0.4pt}\\[0.5cm]
{\large\textbf{End of Document}}
\end{center}
